\scanpage{166}
% position: top sentence
Pour définir un hom \(F\), il faut définir
% position: large central diagram
\diagram{\begin{tikzcd}[column sep=large,row sep=large]
&(E_1\otimes L_2)\unsure{\oplus}(L_1\otimes E_2)=\Sigma
\arrow[dl,"{\mathrm{id}_1\otimes\alpha_2,\alpha_1\otimes\mathrm{id}_2}"']
\arrow[dr,"{\pi_1\otimes\mathrm{id}_{\unsure{L_2|2}},\mathrm{id}_1\otimes\pi_2}"]&\\
E_1\otimes E_2\arrow[dd,"F_{11}"',"\unsure{\times}"{description,pos=0.04}]\arrow[ddrr,"F_{12}" near end]
&&L_1\otimes L_2\arrow[ddll,"F_{21}" near start]\arrow[dd,"F_{22}"]\\
&&\\
E_1\otimes\unsure{\deleted{M_2}\inserted{E_2}}\arrow[dr,"{\mathrm{id}_1\otimes\psi_2,\psi_1\otimes\mathrm{id}_2}"']
&&M_1\otimes M_2\arrow[dl,"{\varpi_1\otimes\mathrm{id}_{\unsure{1|2}},\mathrm{id}_1\otimes\varpi_2}"]\\
&(E_1\otimes M_2)\times(M_1\otimes E_2)=\Pi&
\end{tikzcd}}
% position: right note beside diagram
\margin{\unsure{ces} 4 \unsure{carrés} commu-\linebreak\unsure{tatifs}}
% position: prose below large diagram first condition
Si on veut que \(F\) induise \deleted{l'identité}
\inserted{un hom d'\unsure{extensions}, donc}
\diagram{\begin{tikzcd}L_1\otimes L_2\arrow[r,dash]&L_1\otimes L_2\end{tikzcd}},
il \unsure{f. et s.}
\deleted{\unsure{voir}} que \(\boxed{F_{22}=0}\),
\unsure{si} on veut \inserted{de plus} que \(\unsure{S_0|F}\) induise l'identité sur \(L_1\otimes L_2\),
il \unsure{f. et s.} que \(\boxed{F_{21}=\alpha_1\otimes\alpha_2}\),
si on veut que \(\unsure{F'|F}\) induise l'identité
sur \(M_1\otimes M_2\), il \unsure{f. et s.} que
\(\boxed{F_{12}=\psi_1\otimes\psi_2}\).
\unsure{Si ces} conditions sont \unsure{réalisées}, alors
\unsure{\deleted{deux}|deux} des 4 carrés \unsure{voulus} sont déjà commutatifs
% position: bottom left diagram introduced by savoir
\unsure{savoir}
\diagram{\begin{tikzcd}[row sep=small,column sep=small]
&\Sigma\arrow[dl,dash]\arrow[dr]&\\
E_1\otimes E_2\arrow[dr,"F_{12}=\psi_1\otimes\psi_2"']&0&L_1\otimes L_2\arrow[dl,"F_{22}=0"]\\
&M_1\otimes M_2&
\end{tikzcd}}
% position: bottom right diagram
et
\diagram{\begin{tikzcd}[row sep=small,column sep=small]
\deleted{E_1\otimes E_2}&&L_1\otimes L_2\arrow[dll,"F_{21}"']\arrow[d,"F_{22}=0"]\\
E_1\otimes E_2\arrow[dr]&&M_1\otimes M_2\arrow[dl]\\
&\unsure{\Pi}&
\end{tikzcd}}
% position: bottom right parenthesis
\margin{(\unsure{les quatre} composés sont nuls)}
% position: foot pencil
\margin{165}
